\begin{thebibliography}{52}
\bibitem{r1} Scott Aaronson and Yaoyun Shi. Quantum Lower Bounds for the Collision and the Element Distinctness Problems. J. ACM, 51(4):595–605, July 2004.
\bibitem{r2} Amir Abboud, Ryan Williams, and Huacheng Yu. More Applications of the Polynomial Method to Algorithm Design. In Proceedings of the Twenty-sixth Annual ACM-SIAM Symposium on Discrete Algorithms, SODA ’15, pages 218–230, 2015.
\bibitem{r3} Pankaj K. Agarwal, Herbert Edelsbrunner, Otfried Schwarzkopf, and Emo Welzl. Euclidean Minimum Spanning Trees and Bichromatic Closest Pairs. Discrete \& Computational Geometry, 6(3):407–422, September 1991.
\bibitem{r4} A. Ambainis. Quantum Search Algorithms. SIGACT News, 35(2):22–35, June 2004. doi: 10.1145/992287.992296.
\bibitem{r5} Andris Ambainis. Quantum Walk Algorithm for Element Distinctness. SIAM Journal on Computing, 37(1):210–239, 2007.
\bibitem{r6} Sanjeev Arora and Boaz Barak. Computational Complexity: A Modern Approach. Cambridge University Press, 2009.
\bibitem{r7} Arturs Backurs, Piotr Indyk, and Ludwig Schmidt. On the Fine-Grained Complexity of Empirical Risk Minimization: Kernel Methods and Neural Networks. In Advances in Neural Information Processing Systems, pages 4308–4318, 2017.
\bibitem{r8} Marshall Ball, Alon Rosen, Manuel Sabin, and Prashant Nalini Vasudevan. Proofs of Useful Work. IACR Cryptology ePrint Archive, 2017:203, 2017.
\bibitem{r9} C. Bennett, E. Bernstein, G. Brassard, and U. Vazirani. Strengths and Weaknesses of Quantum Computing. SIAM Journal on Computing, 26(5):1510–1523, 1997.
\bibitem{r10} Jon Louis Bentley and Michael Ian Shamos. Divide-and-Conquer in Multidimensional Space. In Proceedings of the eighth annual ACM symposium on Theory of computing, pages 220–230. ACM, 1976.
\bibitem{r11} Daniel J. Bernstein, Stacey Jeffery, Tanja Lange, and Alexander Meurer. Quantum Algorithms for the Subset-Sum Problem. In Post-Quantum Cryptography, pages 16–33. Springer Berlin Heidelberg, 2013. doi:10.1007/978-3-642-38616-9\_2.
\bibitem{r12} Dominic W. Berry, Andrew M. Childs, Richard Cleve, Robin Kothari, and Rolando D. Somma. Exponential Improvement In Precision for Simulating Sparse Hamiltonians. Forum of Mathematics, Sigma, 5, 2017. doi:10.1017/fms.2017.2.
\bibitem{r13} Sergei N Bespamyatnikh. An Optimal Algorithm for Closest-Pair Maintenance. Discrete \& Computational Geometry, 19(2):175–195, 1998.
\bibitem{r14} Harry Buhrman, Christoph Durr, Mark Heiligman, Peter Hoyer, Frédéric Magniez, Miklos Santha, and Ronald De Wolf. Quantum Algorithms for Element Distinctness. In Proceedings 16th Annual IEEE Conference on Computational Complexity, pages 131–137. IEEE, 2001.
\bibitem{r15} Harry Buhrman, Subhasree Patro, and Florian Speelman. The Quantum Strong ExponentialTime Hypothesis. arXiv preprint, 2019. arXiv:1911.05686.
\bibitem{r16} Timothy M Chan and Ryan Williams. Deterministic APSP, orthogonal vectors, and more: Quickly derandomizing Razborov-Smolensky. In Proceedings of the twenty-seventh annual ACM-SIAM symposium on Discrete algorithms, pages 1246–1255. Society for Industrial and Applied Mathematics, 2016.
\bibitem{r17} Lijie Chen. On the Hardness of Approximate and Exact (Bichromatic) Maximum Inner Product. arXiv preprint, 2018. arXiv:1802.02325.
\bibitem{r18} Lijie Chen and Ryan Williams. An equivalence class for orthogonal vectors. In Proceedings of the Thirtieth Annual ACM-SIAM Symposium on Discrete Algorithms, pages 21–40. SIAM, 2019.
\bibitem{r19} Kenneth L Clarkson. A Randomized Algorithm for Closest-Point Queries. SIAM Journal on Computing, 17(4):830–847, 1988.
\bibitem{r20} Thomas H Cormen, Charles E Leiserson, Ronald L Rivest, and Clifford Stein. Introduction to Algorithms. MIT press, 2009.
\bibitem{r21} Evgeny Dantsin, Vladik Kreinovich, and Alexander Wolpert. On Quantum Versions of Record-breaking Algorithms for SAT. SIGACT News, 36(4):103–108, December 2005. doi: 10.1145/1107523.1107524.
\bibitem{r22} R. David, K. S., and B. Laekhanukit. On the Complexity of Closest Pair via Polar-Pair of Point-Sets. SIAM Journal on Discrete Mathematics, 33(1):509–527, 2019.
\bibitem{r23} Christoph Durr and Peter Hoyer. A quantum algorithm for finding the minimum. arXiv preprint, 1996. arXiv:quant-ph/9607014.
\bibitem{r24} Lov K. Grover. A Fast Quantum Mechanical Algorithm for Database Search. In Proceedings of the Twenty-eighth ACM Symposium on Theory of Computing - STOC '96. ACM Press, 1996. doi:10.1145/237814.237866.
\bibitem{r25} Timon Hertli. Improved Exponential Algorithms for SAT and ClSP. PhD thesis, ETH Zurich, 2015.
\bibitem{r26} Russell Impagliazzo and Ramamohan Paturi. On the Complexity of k-SAT. J. Comput. Syst. Sci., 62(2):367–375, March 2001. doi:10.1006/jcss.2000.1727.
\bibitem{r27} Toshiya Itoh, Tatsuya Nagatani, and Jun Tarui. Explicit Construction for k-Wise Nearly Random Permutations by Iterated Feistel Transform. Workshop on Randomness and Computation, 2005.
\bibitem{r28} Stacey Jeffery. Frameworks for Quantum Algorithms. PhD thesis, University of Waterloo, 2014.
\bibitem{r29} CS Karthik and Pasin Manurangsi. On Closest Pair in Euclidean Metric: Monochromatic is as Hard as Bichromatic. 10th Innovations in Theoretical Computer Science, 2019.
\bibitem{r30} S. Khuller and Y. Matias. A Simple Randomized Sieve Algorithm for the Closest-Pair Problem. Information and Computation, 118(1):34–37, 1995.
\bibitem{r31} Victor Klee. On the Complexity of d-Dimensional Voronoi Diagrams. Archiv der Mathematik, 34(1):75–80, 1980. doi:10.1007/bf01224932.
\bibitem{r32} Jon Kleinberg and Eva Tardos. Algorithm Design. Pearson Education India, 2006.
\bibitem{r33} Frédéric Magniez, Ashwin Nayak, Jérémie Roland, and Miklos Santha. Search via Quantum Walk. SIAM Journal on Computing, 40(1):142–164, 2011.
\bibitem{r34} Chris Marriott and John Watrous. Quantum Arthur–Merlin games. Computational Complexity, 14(2):122–152, 2005. doi:10.1007/s00037-005-0194-x.
\bibitem{r35} Abdullah Mueen, Eamonn Keogh, Qiang Zhu, Sydney Cash, and Brandon Westover. Exact Discovery of Time Series Motifs. In Proceedings of the 2009 SIAM international conference on data mining, pages 473–484. SIAM, 2009.
\bibitem{r36} Alexandros Nanopoulos, Yannis Theodoridis, and Yannis Manolopoulos. C2P: Clustering based on Closest Pairs. In VLDB, 2001.
\bibitem{r37} Michael A Nielsen and Isaac Chuang. Quantum Computation and Quantum Information, 2002.
\bibitem{r38} Ramamohan Paturi and Pavel Pudlák. On the Complexity of Circuit Satisfiability. In Proceedings of the forty-second ACM symposium on Theory of computing, pages 241–250. ACM, 2010.
\bibitem{r39} Ramamohan Paturi, Pavel Pudlák, Michael E Saks, and Francis Zane. An Improved Exponential-Time Algorithm for k-SAT. Journal of the ACM (JACM), 52(3):337–364, 2005.
\bibitem{r40} Michael O Rabin. Probabilistic Algorithms Algorithms and Complexity: New Directions and Recent Results, 1976.
\bibitem{r41} Kunihiko Sadakane, Norito Sugawara, and Takeshi Tokuyama. Quantum Algorithms for Intersection and Proximity Problems. In Peter Eades and Tadao Takaoka, editors, Algorithms and Computation, pages 148–159, Berlin, Heidelberg, 2001. Springer Berlin Heidelberg.
\bibitem{r42} Dominik Scheder and John P Steinberger. PPSZ for General k-SAT-Making Hertli’s Analysis Simpler and 3-SAT Faster. In 32nd Computational Complexity Conference (CCC 2017). Schloss Dagstuhl-Leibniz-Zentrum fuer Informatik, 2017.
\bibitem{r43} T Schoning. A Probabilistic Algorithm for k-SAT and Constraint Satisfaction Problems. In 40th Annual Symposium on Foundations of Computer Science (Cat. No. 99CB37039), pages 410–414. IEEE, 1999.
\bibitem{r44} M. I. Shamos and D. Hoey. Closest-Point Problems. In 16th Annual Symposium on Foundations of Computer Science (sfcs 1975), pages 151–162, 1975.
\bibitem{r45} Mario Szegedy. Quantum Speed-Up of Markov Chain Based Algorithms. In 45th Annual IEEE symposium on foundations of computer science, pages 32–41. IEEE, 2004.
\bibitem{r46} Virginia Vassilevska Williams. Hardness of Easy Problems: Basing Hardness on Popular Conjectures such as the Strong Exponential Time Hypothesis (invited talk). In 10th International Symposium on Parameterized and Exact Computation (IPEC 2015). Schloss Dagstuhl-LeibnizZentrum fuer Informatik, 2015.
\bibitem{r47} Nilton Volpato and Arnaldo Moura. A fast quantum algorithm for the closest bichromatic pair problem, 2010.
\bibitem{r48} Ryan Williams. Pairwise comparison of bit vectors. Theoretical Computer Science Stack Exchange. URL: https://cstheory.stackexchange.com/q/37369.
\bibitem{r49} Ryan Williams. A New Algorithm for Optimal 2-Constraint Satisfaction and Its Implications. Theoretical Computer Science, 348(2-3):357–365, 2005.
\bibitem{r50} Ryan Williams and Huacheng Yu. Finding Orthogonal Vectors In Discrete Structures. In Proceedings of the twenty-fifth annual ACM-SIAM symposium on Discrete algorithms, pages 1867–1877. SIAM, 2014.
\bibitem{r51} Raymond Chi-Wing Wong, Yufei Tao, Ada Wai-Chee Fu, and Xiaokui Xiao. On Efficient Spatial Matching. In Proceedings of the 33rd International Conference on Very Large Data Bases, VLDB ’07, pages 579–590, 2007.
\bibitem{r52} A. C. . Yao. Lower Bounds for Algebraic Computation Trees with Integer Inputs. In 30th Annual Symposium on Foundations of Computer Science, pages 308–313, 1989.
\end{thebibliography}
